\makeatletter
\def\input@path{{./}{../}{../../}{../../../}{../../../../}{preamble/}{../preamble/}{../../preamble/}{../../../preamble/}}
\makeatother

\input{preamble_exams.tex}
\usepackage{needspace}
\usepackage{titling}
\setlength{\droptitle}{-2em}
\pretitle{\par\begingroup\centering\LARGE}
\posttitle{\par\endgroup}
\preauthor{}\postauthor{}
\predate{}\postdate{}
\maketitlehookd{\vspace{-1em}}
\linespread{1.12}
\setlength{\parskip}{5pt}
\newgeometry{margin=0.65in,top=0.55in,bottom=0.65in}

\title{Grade 11 Term 2 Home Practice -- C6\\\large Planning sample size\\[0.25em]
	\normalsize Global Economics -- 2026--2027}
\author{}
\date{}

\begin{document}
	\maketitle

	\setcounter{tocdepth}{1}
	\setcounter{secnumdepth}{2}
	\phantomsection
	\label{toc}
	\tableofcontents

	\section*{Evaluated criteria}
	\begin{enumerate}[label=\textbf{C\arabic*},start=6]
		\item Determines the sample size required to estimate a population's mean or proportion with a specified margin of error, using it to plan economic surveys or financial studies with predetermined precision.
	\end{enumerate}

	\textbf{How to practise.} Try each task before opening its blue solution link. The arrows return to the contents, problem list or current problem.

	\section{Worked Solution Examples}
	\label{sec:worked-solution-examples}

	\SectionProblemsTableOfContents{
		\SectionProblemLink{prob:c6-1}{Planning a port inspection study (C6)} \\
		\SectionProblemLink{prob:c6-2}{Planning a household energy survey (C6)}
	}

	% ----------------------------------------------------------------
	\subsection{Planning a port inspection study (C6)}
	\label{prob:c6-1}

	A logistics agency will estimate the mean cargo clearance time. For planning, use a known population standard deviation of 12 minutes and a 95\% normal critical value of $z^*=1.96$. The agency can sample from a large list of independent eligible cargo records for which a normal approximation is suitable. Its target margin of error is at most 3 minutes.

	\QuestionsTasks[sec:worked-solution-examples]{
		{6}/{sol:c6-1-a}/
		{Start with $E=z^*\sigma/\sqrt n$. Solve the precision inequality for $n$, then give the minimum whole number of records to sample.},
		{6}/{sol:c6-1-b}/
		{Calculate the achieved margin of error for 62 records and for one fewer record. Explain which meets the agency's goal.}
	}

	\clearpage

	\Needspace{5\baselineskip}
	\phantomsection
	\label{sol:c6-1-a}
	\subsubsection*{C6 -- Find the minimum number of records}

	Require $1.96(12)/\sqrt n\le 3$. All quantities are positive, so rearranging gives
	\[
	\sqrt n\ge \frac{1.96(12)}{3}=7.84,
	\qquad n\ge 7.84^2=61.4656.
	\]
	The smallest whole number at or above this bound is 62 records. The agency needs 62 usable measurements under this design.

	\secInicio[prob:c6-1][sec:worked-solution-examples]

	\Needspace{5\baselineskip}
	\phantomsection
	\label{sol:c6-1-b}
	\subsubsection*{C6 -- Verify the rounding decision}

	\[
	E(62)=\frac{1.96(12)}{\sqrt{62}}\approx 2.987\text{ min},
	\qquad E(61)=\frac{1.96(12)}{\sqrt{61}}\approx 3.011\text{ min}.
	\]
	Only 62 meets the 3-minute cap. Since $E(n)$ decreases as $n$ increases, every smaller sample also misses the cap.

	\secInicio[prob:c6-1][sec:worked-solution-examples]

	\clearpage

	% ----------------------------------------------------------------
	\subsection{Planning a household energy survey (C6)}
	\label{prob:c6-2}

	An energy office will estimate the share of households in a city that qualify for a bill subsidy. It will draw a simple random sample without replacement from a complete registry of 20,000 households. It wants a 95\% interval with margin of error at most 4 percentage points. Use $z^*=1.96$ and the planning approximation
	\[
	E\approx 1.96\sqrt{\frac{p_0(1-p_0)}{n}},
	\]
	where $p_0$ is the share used for planning. A small pilot found a share near 0.20.

	\QuestionsTasks[sec:worked-solution-examples]{
		{6}/{sol:c6-2-a}/
		{Convert the precision target to a proportion, solve for the sample size with $p_0=0.20$, and check the expected counts and sampling fraction.},
		{6}/{sol:c6-2-b}/
		{Find a conservative sample size with $p_0=0.50$. Explain why this choice is conservative and why it requires more households.}
	}

	\Needspace{5\baselineskip}
	\phantomsection
	\label{sol:c6-2-a}
	\subsubsection*{C6 -- Use the pilot to calculate a sample size}

	Four percentage points give $E_{\rm target}=0.04$. Squaring the inequality and rearranging gives
	\[
	n\ge \frac{1.96^2(0.20)(0.80)}{0.04^2}=384.16.
	\]
	Plan for at least 385 usable households if the pilot share is suitable. At $p_0=0.20$ the expected counts are 77 qualifying and 308 other households, both above 10. The sampling fraction is $385/20{,}000=1.925\%$. After sampling, check the actual counts and margin.

	\secInicio[prob:c6-2][sec:worked-solution-examples]

	\Needspace{5\baselineskip}
	\phantomsection
	\label{sol:c6-2-b}
	\subsubsection*{C6 -- Plan without a dependable pilot share}

	The product $p_0(1-p_0)=0.25-(p_0-0.50)^2$ is largest at 0.50, so
	\[
	n\ge\frac{1.96^2(0.25)}{0.04^2}=600.25\quad\Longrightarrow\quad n_{\min}=601.
	\]
	This is about 3.0\% of the roster. The plan needs $601-385=216$ extra households because the assumed variability is greater. The 601 records cap the usual normal-formula margin for any observed share, provided the interval's conditions hold. Check both outcome counts after collecting data.

	\secInicio[prob:c6-2][sec:worked-solution-examples]


\end{document}
